% ============================================================
%  Chapter 1 -- The Complex Field, 1.24-1.35
% ============================================================
\rsection{The Complex Field}

\begin{signpost}[Equipment, with one exception]
Nothing in this section is about $\R$, and nothing in it is hard. Rudin is
stocking the shelves: Chapters 3, 7, 8 and 11 all work with complex sequences,
complex series and complex-valued functions, and he wants $\C$ available with
its arithmetic and its absolute value already proved.

\smallskip
The section has a definite shape. Definitions 1.24--1.29 build $\C$ and then
recover the familiar notation $a+bi$; 1.30--1.33 introduce the conjugate and the
absolute value and establish the triangle inequality; 1.34--1.35 prove the
Schwarz inequality. Only two things here have downstream weight: \textbf{1.33(e)},
the triangle inequality, which is used on almost every page of Chapters 3--7,
and \textbf{1.35}, the Schwarz inequality, which is the engine behind the metric
on $\R^k$ in the next section and reappears for integrals in Exercise 6.10 and
for $\mathscr{L}^2$ functions in 11.35.

\smallskip
Watch one point of method. Rudin defines a complex number as an \emph{ordered
pair}, verifies the field axioms by direct computation, and only then --- at
1.27 --- introduces $i$. The mysterious square root of $-1$ is never assumed; it
is produced. That ordering is the whole reason the construction is honest.
\end{signpost}

\ritem{1.24}{Definition}
A \emph{complex number} is an ordered pair $(a,b)$ of real numbers. ``Ordered''
means that $(a,b)$ and $(b,a)$ are regarded as distinct if $a \ne b$.

Let $x = (a,b)$, $y = (c,d)$ be two complex numbers. We write $x = y$ if and
only if $a = c$ and $b = d$. (Note that this definition is not entirely
superfluous; think of equality of rational numbers, represented as quotients of
integers.) We define
\begin{align*}
  x + y &= (a+c,\; b+d)\\
  xy &= (ac-bd,\; ad+bc).
\end{align*}

\begin{unpack}[Where the multiplication rule comes from]
The addition rule is the obvious one. The multiplication rule is not, and Rudin
gives no motivation --- but it is forced, not chosen.

Suppose you want $(a,b)$ to behave like $a + bi$ with $i^2 = -1$, while keeping
the field axioms. Then
\[ (a+bi)(c+di) = ac + adi + bci + bd\,i^2 = (ac - bd) + (ad+bc)i, \]
which is exactly the displayed formula. So the definition is the multiplication
$\C$ \emph{must} have if $i$ is to exist at all.

Rudin declines to say this here, because saying it would presuppose $i$ --- the
very thing he refuses to assume until 1.27. The formula is therefore written
down bare, checked to satisfy the axioms in 1.25, and only afterwards revealed
to mean what you expected. Read 1.24--1.29 as a single argument with the
motivation deliberately withheld until the end.
\end{unpack}

\ritem{1.25}{Theorem}
\emph{These definitions of addition and multiplication turn the set of all
complex numbers into a field, with $(0,0)$ and $(1,0)$ in the role of $0$ and
$1$.}

\begin{proof}
We simply verify the field axioms, as listed in Definition 1.12. (Of course, we
use the field structure of $\R$.) Let $x = (a,b)$, $y = (c,d)$, $z = (e,f)$.

\smallskip
(A1) is clear.

(A2) $x+y = (a+c,\,b+d) = (c+a,\,d+b) = y+x$.

(A3) \begin{align*}
  (x+y)+z &= (a+c,\,b+d)+(e,f)\\
          &= (a+c+e,\; b+d+f)\\
          &= (a,b)+(c+e,\,d+f) = x+(y+z).
\end{align*}

(A4) $x + 0 = (a,b)+(0,0) = (a,b) = x$.

(A5) Put $-x = (-a,-b)$. Then $x + (-x) = (0,0) = 0$.

\smallskip
(M1) is clear.

(M2) $xy = (ac-bd,\,ad+bc) = (ca-db,\,da+cb) = yx$.

(M3) \begin{align*}
  (xy)z &= (ac-bd,\,ad+bc)(e,f)\\
        &= (ace-bde-adf-bcf,\; acf-bdf+ade+bce)\\
        &= (a,b)(ce-df,\,cf+de) = x(yz).
\end{align*}

(M4) $1x = (1,0)(a,b) = (a,b) = x$.

(M5) If $x \ne 0$ then $(a,b) \ne (0,0)$, which means that at least one of the
real numbers $a$, $b$ is different from $0$. Hence $a^2+b^2 > 0$, by Proposition
1.18(d), and
we can define\kn{Exactly the leverage promised at 1.18(d). Squares are
nonnegative and not both $a,b$ vanish, so $a^2+b^2>0$ --- and therefore the
divisions below are legal. The construction of $\C$ depends on an order
property of $\R$, which is worth noticing given that $\C$ itself cannot be
ordered (Exercise 1.8).}
\[ \frac1x = \left(\frac{a}{a^2+b^2},\; \frac{-b}{a^2+b^2}\right). \]
Then
\[ x \cdot \frac1x = (a,b)\left(\frac{a}{a^2+b^2},\;
   \frac{-b}{a^2+b^2}\right) = (1,0) = 1. \]

\smallskip
(D) \begin{align*}
  x(y+z) &= (a,b)(c+e,\,d+f)\\
         &= (ac+ae-bd-bf,\; ad+af+bc+be)\\
         &= (ac-bd,\,ad+bc)+(ae-bf,\,af+be)\\
         &= xy + xz. \qedhere
\end{align*}
\end{proof}

\begin{deeper}[Nothing here is deep, and that is the point]
Every line is arithmetic in $\R$. That is precisely what makes the theorem
worth stating: it shows the existence of $\C$ requires no new assumptions
whatever --- given $\R$, the complex field comes free, as a set of pairs.

Compare the situation for $\R$ itself, where Theorem 1.19 needed a page of
Dedekind cuts. The difference is that $\C$ adds no \emph{completeness}; it only
adds algebra. That algebra turns out to be enormously powerful --- 8.8 will show
$\C$ is algebraically closed --- but it costs nothing to obtain.
\end{deeper}

\ritem{1.26}{Theorem}
\emph{For any real numbers $a$ and $b$ we have}
\[ (a,0)+(b,0) = (a+b,0), \qquad (a,0)(b,0) = (ab,0). \]

The proof is trivial.

Theorem 1.26 shows that the complex numbers of the form $(a,0)$ have the same
arithmetic properties as the corresponding real numbers $a$. We can therefore
identify $(a,0)$ with $a$. This identification gives us the real field as a
subfield of the complex field.\kd{The same manoeuvre as Step 9 of the appendix,
where $\Q$ is identified with the rational cuts $\Q^*$. In both cases a
\emph{copy} of the smaller field sits inside the larger one, arithmetic and
order preserved, and we agree to stop distinguishing the copy from the original.
Rudin explicitly links the two identifications in Step 9; when you read the
appendix, come back to this paragraph.}

The reader may have noticed that we defined the complex numbers without any
reference to the mysterious square root of $-1$. We now show that the notation
$(a,b)$ is equivalent to the more customary $a+bi$.

\ritem{1.27}{Definition}
$i = (0,1)$.

\ritem{1.28}{Theorem}
$i^2 = -1$.

\begin{proof}
$i^2 = (0,1)(0,1) = (-1,0) = -1$.
\end{proof}

\ritem{1.29}{Theorem}
\emph{If $a$ and $b$ are real, then $(a,b) = a+bi$.}

\begin{proof}
\begin{align*}
  a + bi &= (a,0)+(b,0)(0,1)\\
         &= (a,0)+(0,b) = (a,b). \qedhere
\end{align*}
\end{proof}

\begin{unpack}[What 1.27--1.29 actually accomplish]
Three lines of computation, and they retire a piece of folklore. The usual story
is that $i$ is postulated --- an imaginary object adjoined by fiat because we
wish $x^2+1$ had a root. Rudin's construction assumes no such thing. $\C$ is
built from ordered pairs of reals, its arithmetic is defined outright, and then
one particular pair, $(0,1)$, is observed to square to $-1$.

So $i$ is not an assumption; it is a \emph{discovery} about a field that already
existed. And 1.29 shows the pair notation and the $a+bi$ notation describe the
same object, so you may use whichever is convenient --- pairs when checking
axioms, $a+bi$ when computing.
\end{unpack}

\ritem{1.30}{Definition}
If $a$, $b$ are real and $z = a+bi$, then the complex number $\bar z = a-bi$ is
called the \emph{conjugate} of $z$. The numbers $a$ and $b$ are the \emph{real
part} and the \emph{imaginary part} of $z$, respectively.

We shall occasionally write
\[ a = \operatorname{Re}(z), \qquad b = \operatorname{Im}(z). \]

\begin{center}
\begin{tikzpicture}[x=1mm, y=1mm]
  \draw[axis] (-16,0) -- (34,0);  \node[lbl, anchor=west] at (34,0) {$\operatorname{Re}$};
  \draw[axis] (0,-20) -- (0,22);  \node[lbl, anchor=south] at (0,22) {$\operatorname{Im}$};
  % z
  \draw[vec] (0,0) -- (22,14);
  \node[ptfill] at (22,14) {};
  \node[mlbl, anchor=south west] at (22,14) {$z=(a,b)=a+bi$};
  \node[mlbl, anchor=south east] at (11,9) {$|z|$};
  % conjugate
  \draw[vec] (0,0) -- (22,-14);
  \node[ptfill] at (22,-14) {};
  \node[mlbl, anchor=north west] at (22,-14) {$\bar z=a-bi$};
  \draw[guide] (22,14) -- (22,-14);
  % components
  \draw[guide] (0,14) -- (22,14);
  \node[mlbl, anchor=east] at (-1,14) {$b$};
  \node[mlbl, anchor=north] at (22,-0.5) {$a$};
  \node[ptfill, minimum size=2.6pt] at (22,0) {};
  % i
  \node[ptfill, minimum size=2.6pt] at (0,8) {};
  \node[mlbl, anchor=east] at (-1,8) {$i=(0,1)$};
\end{tikzpicture}
\end{center}
\figcap{Conjugation is reflection in the real axis, and
$|z|=\sqrt{a^2+b^2}$ is distance from the origin. Neither is a new assumption
--- both are read off Definition 1.24.}

\ritem{1.31}{Theorem}
\emph{If $z$ and $w$ are complex, then}
\begin{enumerate}[label=(\alph*), leftmargin=2.2em, itemsep=1pt, topsep=3pt]
  \item \emph{$\overline{z+w} = \bar z + \bar w$,}
  \item \emph{$\overline{zw} = \bar z \cdot \bar w$,}
  \item \emph{$z + \bar z = 2\operatorname{Re}(z)$,
        $z - \bar z = 2i\operatorname{Im}(z)$,}
  \item \emph{$z\bar z$ is real and positive (except when $z = 0$).}
\end{enumerate}

\begin{proof}
(a), (b), and (c) are quite trivial. To prove (d), write $z = a+bi$, and note
that $z\bar z = a^2+b^2$.
\end{proof}

\begin{deeper}[Conjugation is a field automorphism]
Parts (a) and (b) say that $z \mapsto \bar z$ preserves both operations, and
since $\bar{\bar z} = z$ it is a bijection. In the language of algebra,
conjugation is an automorphism of $\C$ fixing $\R$ pointwise --- and it is the
only nontrivial one.

The practical use is a symmetry principle you will meet repeatedly: any
polynomial identity with real coefficients that holds for $z$ also holds for
$\bar z$. That is why non-real roots of real polynomials come in conjugate
pairs, and why 8.8 (the algebraic completeness of $\C$) has the consequences it
does for real polynomials.

Part (d) is the one with immediate work to do: it is what makes Definition 1.32
possible, since $z\bar z \ge 0$ is exactly the condition needed before taking a
square root by 1.21.
\end{deeper}

\ritem{1.32}{Definition}
If $z$ is a complex number, its \emph{absolute value} $|z|$ is the non-negative
square root of $z\bar z$; that is, $|z| = (z\bar z)^{1/2}$.

The existence (and uniqueness) of $|z|$ follows from Theorem 1.21 and part (d)
of Theorem 1.31.\kn{Both citations are load-bearing, and note the gap they leave.
For $z \ne 0$, 1.31(d) makes $z\bar z$ a positive real and 1.21 --- which
requires a positive input --- returns its unique positive square root. For
$z = 0$ we have $z\bar z = 0$, which 1.21 does not cover; that case is settled
separately by the convention $|0| = 0$, recorded in 1.33(a). This is the first
place the work of \S1.21 gets spent.}

Note that when $x$ is real, then $\bar x = x$, hence $|x| = \sqrt{x^2}$. Thus
$|x| = x$ if $x \ge 0$, $|x| = -x$ if $x < 0$.

\ritem{1.33}{Theorem}
\emph{Let $z$ and $w$ be complex numbers. Then}
\begin{enumerate}[label=(\alph*), leftmargin=2.2em, itemsep=1pt, topsep=3pt]
  \item \emph{$|z| > 0$ unless $z = 0$, $|0| = 0$,}
  \item \emph{$|\bar z| = |z|$,}
  \item \emph{$|zw| = |z|\,|w|$,}
  \item \emph{$|\operatorname{Re}(z)| \le |z|$,}
  \item \emph{$|z+w| \le |z|+|w|$.}
\end{enumerate}

\begin{proof}
(a) and (b) are trivial. Put $z = a+bi$, $w = c+di$, with $a,b,c,d$ real. Then
\[ |zw|^2 = (ac-bd)^2+(ad+bc)^2 = (a^2+b^2)(c^2+d^2) = |z|^2|w|^2 \]
or $|zw|^2 = (|z|\,|w|)^2$. Now (c) follows from the uniqueness assertion of
Theorem 1.21.

To prove (d), note that $a^2 \le a^2+b^2$, hence
\[ |a| = \sqrt{a^2} \le \sqrt{a^2+b^2}. \]

To prove (e), note that $\bar z w$ is the conjugate of $z\bar w$, so that
$z\bar w + \bar z w = 2\operatorname{Re}(z\bar w)$. Hence
\begin{align*}
  |z+w|^2 &= (z+w)(\overline{z+w}) = z\bar z + z\bar w + \bar z w + w\bar w\\
          &= |z|^2 + 2\operatorname{Re}(z\bar w) + |w|^2\\
          &\le |z|^2 + 2|z\bar w| + |w|^2\\
          &= |z|^2 + 2|z|\,|w| + |w|^2 = (|z|+|w|)^2.
\end{align*}
Now (e) follows by taking square roots.
\end{proof}

\begin{unpack}[The standard proof of a triangle inequality]
The shape of this argument recurs throughout analysis, and it is worth extracting
because you will need to reproduce it: to bound $|z+w|$, do not work with
$|z+w|$ --- work with $|z+w|^2$, expand it as $(z+w)(\overline{z+w})$, bound the
cross-term, and take square roots at the end.

The cross-term is the only real step, and it goes in two moves: 1.31(c) turns
$z\bar w + \bar z w$ into $2\operatorname{Re}(z\bar w)$, then (d) bounds
$\operatorname{Re}$ by modulus and (c) factors it as $|z||w|$.

The identical strategy proves 1.37(e) for $\R^k$ in the next section, with the
Schwarz inequality playing the part of (d). Recognising the pattern once saves
reading the second proof.
\end{unpack}

\begin{center}
\begin{tikzpicture}[x=1mm, y=1mm]
  \draw[seg] (0,0) -- (26,5);
  \draw[seg] (26,5) -- (34,24);
  \draw[thin] (0,0) -- (34,24);
  \node[ptfill, minimum size=2.4pt] at (0,0) {};
  \node[ptfill, minimum size=2.4pt] at (26,5) {};
  \node[ptfill, minimum size=2.4pt] at (34,24) {};
  \node[mlbl, anchor=north] at (13,-1.2) {$|z|$};
  \node[mlbl, anchor=west] at (32,13.4) {$|w|$};
  \node[mlbl, anchor=south east] at (13,12.2) {$|z{+}w|$};
  \node[mlbl, anchor=north] at (26,2.8) {$z$};
  \node[mlbl, anchor=south] at (34,25.6) {$z{+}w$};
  \node[mlbl, anchor=east] at (-1.6,0) {$0$};
\end{tikzpicture}
\end{center}
\figcap{Part (e) is about an actual triangle: the straight route from $0$ to
$z+w$ never beats the detour through $z$. Equality is the degenerate triangle
--- $z$ and $w$ pointing the same way, which is exactly
$\operatorname{Re}(z\bar w) = |z\bar w|$, the one inequality in the proof.}

\begin{deeper}[Multiplication rotates and scales --- (c) drawn instead of computed]
Part (c) also has a picture. Multiplying every point of the plane by $w$
rotates it by the angle of $w$ and stretches it by the factor $|w|$: the
triangle with vertices $0, w, zw$ is the triangle $0, 1, z$ so rotated and
stretched.

\medskip
\begin{center}
\begin{tikzpicture}[x=1mm, y=1mm]
  \draw[axis] (-4,0) -- (26,0);
  \draw[axis] (0,-2) -- (0,30);
  \fill[black!12] (0,0) -- (9,0) -- (18,9) -- cycle;
  \draw[thin] (0,0) -- (9,0) -- (18,9) -- cycle;
  \fill[pattern=north east lines, pattern color=black!45]
    (0,0) -- (9,9) -- (9,27) -- cycle;
  \draw[thin] (0,0) -- (9,9) -- (9,27) -- cycle;
  \node[ptfill, minimum size=2.4pt] at (9,0) {};
  \node[ptfill, minimum size=2.4pt] at (18,9) {};
  \node[ptfill, minimum size=2.4pt] at (9,9) {};
  \node[ptfill, minimum size=2.4pt] at (9,27) {};
  \node[mlbl, anchor=north] at (9,-1.6) {$1$};
  \node[mlbl, anchor=west] at (19.6,9) {$z$};
  \node[mlbl, anchor=south west] at (11,7.4) {$w$};
  \node[mlbl, anchor=west] at (10.6,27) {$zw$};
  \node[mlbl, anchor=north east] at (-1,-0.4) {$0$};
\end{tikzpicture}
\end{center}
\figcap{The shaded triangle $0,1,z$ and the hatched triangle $0,w,zw$ are
similar, with ratio $|w|$. (Drawn for $z = 2+i$, $w = 1+i$; the similarity is
general.)}

\smallskip
Part (c) is the stretch factor read off the picture: the side $0 \to zw$ is
the side $0 \to z$ magnified by $|w|$, so $|zw| = |z||w|$. Rudin proves it by
the $z\bar z$ computation instead, because the rotation--scaling description
presupposes angles --- sines and cosines --- which do not officially exist
until Chapter 8. The picture is why the computation was never in doubt; the
computation is what makes the picture legal. Taking $w = \bar z$ folds the
rotation flat onto the real axis: $z\bar z = |z|^2$ is this figure with the
hatched triangle lying on its back.
\end{deeper}

\ritem{1.34}{Notation}
If $x_1, \dots, x_n$ are complex numbers, we write
\[ x_1 + x_2 + \cdots + x_n = \sum_{j=1}^n x_j. \]

We conclude this section with an important inequality, usually known as the
\emph{Schwarz inequality}.

\ritem{1.35}{Theorem}
\emph{If $a_1,\dots,a_n$ and $b_1,\dots,b_n$ are complex numbers, then}
\[ \left|\sum_{j=1}^n a_j\bar b_j\right|^2 \le
   \sum_{j=1}^n |a_j|^2 \sum_{j=1}^n |b_j|^2. \]

\begin{proof}
Put $A = \Sigma|a_j|^2$, $B = \Sigma|b_j|^2$, $C = \Sigma a_j\bar b_j$ (in all
sums in this proof, $j$ runs over the values $1,\dots,n$). If $B = 0$, then
$b_1 = \cdots = b_n = 0$, and the conclusion is trivial. Assume therefore that
$B > 0$. By Theorem 1.31 we have
\begin{align*}
  \sum|Ba_j - Cb_j|^2
    &= \sum(Ba_j-Cb_j)(\overline{Ba_j-Cb_j})\\
    &= B^2\sum|a_j|^2 - B\bar C\sum a_j\bar b_j\\
    &\qquad {}- BC\sum \bar a_j b_j + |C|^2\sum|b_j|^2\\
    &= B^2A - B|C|^2\\
    &= B(AB-|C|^2).
\end{align*}
Since each term in the first sum is nonnegative, we see that
\[ B(AB - |C|^2) \ge 0. \]
Since $B > 0$, it follows that $AB - |C|^2 \ge 0$. This is the desired
inequality.
\end{proof}

\begin{center}
\begin{tikzpicture}[x=1mm, y=1mm]
  \node[rmbox, text width=30mm] (sq) at (0,20)
    {$\sum|Ba_j-Cb_j|^2 \ge 0$\\ \textit{(squares are $\ge 0$)}};
  \node[rmbox, text width=30mm] (ex) at (62,20)
    {expand: it equals\\$B(AB-|C|^2)$};
  \draw[rmarrow] (sq) -- node[lbl, anchor=south, inner sep=2.5pt]
    {1.31, algebra} (ex);
  \node[rmbox, text width=30mm] (dv) at (62,0)
    {divide by $B>0$\\($B=0$: trivial case)};
  \draw[rmarrow] (ex) -- (dv);
  \node[rmbox, text width=24mm] (cc) at (0,0) {$|C|^2 \le AB$};
  \draw[rmarrow] (dv) -- (cc);
\end{tikzpicture}
\end{center}
\figcap{The proof as a flow diagram: one nonnegative square, expanded and
collapsed. All the invention is in choosing \emph{which} square --- the box
below recovers that choice.}

\begin{unpack}[Where $Ba_j - Cb_j$ comes from]
This is the second rabbit of the chapter, after the $h$ of 1.21, and like that
one it is reverse-engineered rather than guessed.

The only fact available is that a sum of squared moduli is nonnegative. So the
proof must begin with $\sum|\,\cdot\,|^2 \ge 0$ for some cleverly chosen
expression, and end with the desired inequality. Work backwards: we want
$AB - |C|^2 \ge 0$, so we want the sum to evaluate to a positive multiple of
$AB - |C|^2$.

Try $\sum|\lambda a_j - \mu b_j|^2$ with constants to be chosen. Expanding gives
\[ |\lambda|^2 A - \lambda\bar\mu\, \bar C - \bar\lambda\mu\, C + |\mu|^2 B . \]
Choosing $\lambda = B$ (real) and $\mu = C$ makes the middle terms
$-B\bar C C - BC\bar C = -2B|C|^2$ and the last term $|C|^2B$, so the whole
collapses to $B^2A - B|C|^2 = B(AB - |C|^2)$. The factor $B$ is then stripped
off using $B>0$.

\smallskip
So the choice is not inspired --- it is what you get by writing down the general
quadratic and solving for the constants that make it collapse. The hypothesis
$B > 0$ and the separate treatment of $B = 0$ exist purely because the division
at the last step needs it.
\end{unpack}

\begin{deeper}[Equality, and what the inequality means]
Equality holds precisely when $(a_j)$ and $(b_j)$ are \emph{linearly dependent}.
That is Exercise 1.15, and the proof above shows why --- with a caveat about how
it is usually stated. When $B>0$, equality forces $\sum|Ba_j-Cb_j|^2 = 0$, and a
sum of nonnegative terms vanishes only if every term does, giving
$Ba_j = Cb_j$ for all $j$: the two vectors are proportional. When $B = 0$ every
$b_j$ vanishes, both sides of the inequality are $0$, and equality holds
trivially --- but ``proportional'' now says nothing, since $Ba_j = Cb_j$ reduces
to $0=0$. Linear dependence is the formulation that covers both cases.

Geometrically, once 1.36 defines the inner product on $\R^k$, this says
$|\mathbf{x}\cdot\mathbf{y}| \le |\mathbf{x}||\mathbf{y}|$, which is the
statement that $\cos\theta$ lies between $-1$ and $1$. Rudin proves the
inequality first and lets the geometry follow, rather than the reverse; the
inequality is the substance and the picture is the illustration.

Its reach is long. It gives the triangle inequality in $\R^k$ at 1.37(e), hence
the metric that all of Chapter 2 is built on; it reappears for integrals in
Exercise 6.10 --- H\"older's inequality, of which Rudin there remarks
``Theorem 1.35 is a very special case'' --- and for $\mathscr{L}^2$ functions
at 11.35, where it is the same statement about a different inner product.
\end{deeper}
